\chapter{2002年上海卷（春）}
\section{填空题}
\begin{problem}
函数 \( y = \frac{1}{3 - 2x - x^2} \) 的定义域为\fillinblank{}.
\end{problem}

\begin{problem}
若椭圆的两个焦点坐标为 \(F_{1}(-1,0), F_{2}(5,0)\)，长轴长为 10，则椭圆的方程为\fillinblank{}.
\end{problem}

\begin{problem}
若全集 \(I = \R\)，\(f(x)\)，\(g(x)\) 均为 \(x\) 的二次函数，\(P = \{x \mid f(x) < 0\}\)，\(Q = \{x \mid g(x) \geqslant 0\}\)，则不等式组 \(\begin{cases} f(x) < 0, \\ g(x) < 0 \end{cases}\) 的解集可用 \(P, Q\) 表示为\fillinblank{}.
\end{problem}

\begin{problem}
设 \( f(x) \) 是定义在 \( \R \) 上的奇函数. 若当 \( x \geqslant 0 \) 时，\( f(x) = \log_3(1 + x) \)，则 \( f(-2) = \)\fillinblank{}.
\end{problem}

\begin{problem}
若在 \(\left(\sqrt[5]{x} - \frac{1}{x}\right)^n\) 的展开式中，第 4 项是常数项，则 \(n = \)\fillinblank{}.
\end{problem}

\begin{problem}
已知 \( f(x) = \sqrt{\frac{1 - x}{1 + x}} \)，若 \( \alpha \in \left(\frac{\pi}{2}, \pi\right) \)，则 \( f(\cos \alpha) + f(-\cos \alpha) \) 可化简为\fillinblank{}.
\end{problem}

\begin{problem}
六位身高全不相同的同学拍照留念. 摄影师要求前后两排各三人，则后排每人均比前排同学高的概率是\fillinblank{}.
\end{problem}

\begin{problem}
设曲线 \(C_1\) 和 \(C_2\) 的方程分别为 \(F_1(x,y) = 0\) 和 \(F_2(x,y) = 0\)，则点 \(P(a,b) \notin C_1 \cap C_2\) 的一个充分条件为\fillinblank{}.
\end{problem}

\begin{problem}
若 \( f(x) = 2 \sin \omega x \left(0 < \omega < 1\right) \) 在区间 \( \left[0, \frac{\pi}{3}\right] \) 上的最大值是 \( \sqrt{2} \) ，则 \( \omega = \)\fillinblank{}.
\end{problem}

\begin{problem}
如图表示一个正方体表面的一种展开图，图中的四条线段 \(AB\)，\(CD\)，\(EF\) 和 \(GH\) 在原正方体中相互异面的有 对\fillinblank{}.

\begin{center}
\bitmapinclude[page=1,pagebox=mediabox,trim=99pt 122pt 777pt 405pt,clip,max width=0.42\linewidth,max totalheight=4.6cm]{source/春季高考/2002/2002春季上海.pdf}
\end{center}
\end{problem}

\begin{problem}
如图所示，客轮以速度 \(2v\) 由 \(A\) 至 \(B\) 再到 \(C\) 匀速航行，货轮从 \(AC\) 的中点 \(D\) 出发，以速度 \(v\) 沿直线匀速航行，将货物送达客轮. 已知 \(AB \perp BC\)，且 \(AB = BC = 50\text{海里}\). 若两船同时出发，则两船相遇之处距 \(C\) 点\fillinblank{}\(\text{海里}\). (结果精确到小数点后 1 位)

\begin{center}
\bitmapinclude[page=1,pagebox=mediabox,trim=540pt 717pt 507pt 27pt,clip,max width=0.30\linewidth,max totalheight=3.4cm]{source/春季高考/2002/2002春季上海.pdf}
\end{center}
\end{problem}

\begin{problem}
如图，若从点 \(O\) 所作的两条射线 \(OM, ON\) 分别有点 \(M_{1}, M_{2}\) 与点 \(N_{1}, N_{2}\)，则三角形面积之比 \(\frac{S_{\triangle O M_{1}N_{1}}}{S_{\triangle O M_{2}N_{2}}} = \frac{OM_{1}}{OM_{2}} \cdot \frac{ON_{1}}{ON_{2}}\)，若从点 \(O\) 所作的不在同一平面内的三条射线 \(OP, OQ\) 和 \(OR\) 上，分别有点 \(P_{1}, P_{2}\)，点 \(Q_{1}, Q_{2}\) 和点 \(R_{1}, R_{2}\) 则类似的结论为\fillinblank{}.

\begin{center}
\bitmapinclude[page=1,pagebox=mediabox,trim=423pt 526pt 453pt 216pt,clip,max width=0.64\linewidth,max totalheight=4.6cm]{source/春季高考/2002/2002春季上海.pdf}
\end{center}
\end{problem}

\section{单选题}
\begin{problem}
若 \(a, b, c\) 为任意向量，\(m \in \R\)，则下列等式不一定成立的是（\quad）

\choices
    {\((a + b) + c = a + (b + c)\)}
    {\((a + b) \cdot c = a \cdot c + b \cdot c\)}
    {\(m(a + b) = ma + mb\)}
    {\((a\cdot b)c = a(b\cdot c)\)}
\end{problem}

\begin{problem}
在 \(\triangle ABC\) 中，若 \(2\cos B\sin A = \sin C\)，则 \(\triangle ABC\) 的形状一定是（\quad）

\choices
    {等腰直角三角形}
    {直角三角形}
    {等腰三角形}
    {等边三角形}
\end{problem}

\begin{problem}
设 \(A > 0, a \neq 1\)，函数 \(y = \log_a x\) 的反函数和 \(y = \log_a \frac{1}{x}\) 的反函数的图象关于（\quad）

\choices
    {\(x\) 轴对称}
    {\(y\) 轴对称}
    {\(y = x\) 对称}
    {原点对称}
\end{problem}

\begin{problem}
设 \(\{\alpha_{n}\}\) (\(n \in \mathbf{N}\)) 是等差数列，\(S_{n}\) 是其前 \(n\) 项的和，且 \(S_{5} < S_{6}\)，\(S_{6} = S_{7} > S_{8}\)，则下列结论错误的是（\quad）

\choices
    {\(d < 0\)}
    {\(\alpha_{7} = 0\)}
    {\(S_{9} > S_{5}\)}
    {\(S_{6}\) 和 \(S_{7}\) 均为 \(S_{n}\) 的最大值}
\end{problem}

\section{解答题}
\begin{problem}
已知 \(z\)，\(\omega\) 为复数，\((1 + 3\i)z\) 为纯虚数，\(\omega = \frac{z}{2 + \i}\) 且 \(|\omega| = 5\sqrt{2}\)，求 \(\omega\).
\end{problem}

\begin{problem}
已知 \(F_{1}\)，\(F_{2}\) 为双曲线 \(\frac{x^{2}}{a^{2}} - \frac{y^{2}}{b^{2}} = 1 (a > 0, b > 0)\) 的焦点. 过 \(F_{2}\) 作垂直 \(x\) 轴的直线交双曲线于点 \(P\)，且 \(\angle PF_{1}F_{2} = 30^{\circ}\)，求双曲线的渐近方程.
\end{problem}

\begin{problem}
如图，三棱柱 \(OAB - O_{1}A_{1}B_{1}\)，平面 \(OBB_{1}O_{1} \perp\) 平面 \(OAB\)，\(\angle O_{1}OB = 60^{\circ}\)，\(\angle AOB = 90^{\circ}\)，且 \(OB = OO_{1} = 2\)，\(OA = \sqrt{3}\). 求：
\begin{enumerate}[label={（\arabic*）},itemsep=0.2em]
\item 二面角 \(O_{1} - AB - O\) 的大小；
\item 异面直线 \(A_{1}B\) 与 \(AO_{1}\) 所成角的大小. (上述结果用反三角函数值表示)
\end{enumerate}

\begin{center}
\bitmapinclude[page=1,pagebox=mediabox,trim=846pt 239pt 30pt 473pt,clip,max width=0.45\linewidth,max totalheight=4.8cm]{source/春季高考/2002/2002春季上海.pdf}
\end{center}
\end{problem}

\begin{problem}
已知函数 \( f(x) = a^x + \frac{x - 2}{x + 1} (a > 1) \).
\begin{enumerate}[label={（\arabic*）},itemsep=0.2em]
\item 证明：函数 \( f(x) \) 在 \( (-1, +\infty) \) 上为增函数；
\item 用反证法证明方程 \( f(x)=0 \) 没有负数根.
\end{enumerate}
\end{problem}

\begin{problem}
某公司全年的利润为 \(b\text{元}\)，其中一部分作为奖金发给 \(n\) 位职工，奖金分配方案如下：首先将职工按工作业绩 (工作业绩均不相同) 从大到小，由 1 到 \(n\) 排序. 第 1 位职工得奖金 \(\frac{b}{n}\text{元}\)，然后再将余额除以 \(n\) 发给第 2 位职工，按此方法将奖金逐一发给每位职工，并将最后剩余部分作为公司发展基金.
\begin{enumerate}[label={（\arabic*）},itemsep=0.2em]
\item 设 \( a_k (1 \leqslant k \leqslant m) \) 为第 \( k \) 位职工所得奖金额，试求 \( a_2, a_3 \)，并用 \( k, n \) 和 \( b \) 表示 \( a \) (不必证明)；
\item 证明 \(a_{k} > a_{k + 1}\) (\(k = 1, 2, \dots, n - 1\))，并解释此不等式关于分配原则的实际意义；
\item 发展基金与 \(n\) 和 \(b\) 有关，记为 \(P_{n}(b)\)，对常数 \(b\)，当 \(n\) 变化时，求 \(\lim_{n \to \infty} P_{n}(b)\).
\end{enumerate}
\end{problem}

\begin{problem}
对于函数 \( f(x) \)，若存在 \( x_0 \in \R \)，使 \( f(x_0) = x_0 \) 成立，则称 \( x_0 \) 为 \( f(x) \) 的不动点. 已知函数 \( f(x) = ax^2 + (b + 1)x + (b - 1) (a \neq 0) \).
\begin{enumerate}[label={（\arabic*）},itemsep=0.2em]
\item 当 \(a = 1, b = -2\) 时，求函数 \(f(x)\) 的不动点；
\item 若对任意实数 \(b\)，函数 \(f(x)\) 恒有两个相异的不动点，求 \(a\) 的取值范围；
\item 在 （2） 的条件下，若 \(y = f(x)\) 图上 \(A\)，\(B\) 两点的横坐标是函数 \(f(x)\) 的不动点，且 \(A\)，\(B\) 两点关于直线 \(y = kx + \frac{1}{2a^2 + 1}\) 对称，求 \(b\) 的最小值.
\end{enumerate}
\end{problem}
